\documentstyle[%


righttag%


,11pt%


,amssymb%


,russian%               remove in Eng.


,izvru%                 Set izven% in Eng.


]{amsart}





% \not\to ?





\newcommand{\dom}{\operatorname{dom}}


\newcommand{\epi}{\operatorname{epi}}


\newcommand{\mes}{\operatorname{mes}}


\newcommand{\supp}{\operatorname{supp}}


\newcommand{\tts}[1]{}





\theoremstyle{plain}


\theorembodyfont{\it}


\newtheorem{thms}{\hskip\parindent Теорема}[section]


\newtheorem{lems}{\hskip\parindent Лемма}[section]





\theoremstyle{definition}


\theorembodyfont{\rm}


\newtheorem{cors}{\hskip\parindent Следствие}[section]


\newtheorem{defns}{\hskip\parindent Определение}[section]


\newtheorem{notes}{\hskip\parindent Замечание}[section]





\numberwithin{equation}{section}





%\hoffset=-15mm%                  For Eng.


\setcounter{page}{40}


\god{2005}


\nomer{6~(517)} %                        Remove in Eng.


%\nomer{Vol.\,49, No.\,6}%               Set in Eng.


%\str{\pageref{begin}--\pageref{end}}%  Set in Eng.


\udk{517.977}





\author{\it В.С.\,Гаврилов, М.И.\,Сумин}


% NB:   ^^ remove in Eng.


\title{Параметрическая задача субоптимального\\ управления системой


Гурса--Дарбу\\


с поточечным фазовым ограничением}


\thanks{Работа выполнена при


финансовой поддержке Российского фонда фундаментальных исследований


(проекты \No\No\,01-01-00979, 04-01-00460) и конкурсного центра 


фундаментального естествознания при Санкт-Петербургском государственном 


университете (проект \No\,Е02-1.0-173).}





\date{}


%\pagestyle{myheadings}%         Set in Eng.


\pagestyle{plain} %              Remove in Eng.


\begin{document}


%\thispagestyle{plain}%          Set in Eng.


\maketitle


\label{begin}











\section*{Введение}





Основным предположением теории необходимых условий оптимальности в задачах


оптимального управления традиционно является предположение о существовании


оптимального элемента, каковым может служить как обычное, т.\,е. измеримое по


Лебегу, так и обобщенное управление \cite{gam77}, \cite{warga}. Как


известно \cite{young}, в случае обычного оптимального


управления именно это обстоятельство привело к словосочетанию


``наивная теория оптимального управления''. Заметим,


что для задач оптимального управления обыкновенными


дифференциальными уравнениями, а также 


полулинейными уравнениями в частных производных


существование оптимальных обобщенных \cite{gam77}, \cite{warga}


элементов дается ``практически даром''.


Естественно, постулирование факта существования сопряжено с наложением весьма


жестких условий на исходные данные задач.





В то же время специфика управляемых


распределенных систем такова, что несуществование обычных оптимальных


управлений и одновременно невозможность расширения задачи в каком-либо


естественном смысле (А.Ф.\,Филиппов, Л.\,Янг, Р.В.\,Гамкрелидзе, Дж.\,Варга)


представляет собой их типичное


свойство, которым они существенно отличаются от


управляемых систем обыкновенных дифференциальных уравнений.  Содержательные


примеры таких задач многочисленны и, например, в случае управляемой системы


Гурса--Дарбу могут быть найдены в \cite{SuminD}, \cite{danmatv}. Анализ


подобного примера для задачи оптимизации системы Гурса--Дарбу с поточечным


фазовым ограничением приводится в заключительной части данной


работы. Этот анализ, а также анализ многих других аналогичных примеров


(напр., \cite{SuminD}, \cite{danmatv}) позволяет заключить, что в


общем нелинейном случае системы Гурса--Дарбу ее нельзя расширить в смысле


А.Ф.\,Филиппова, Л.\,Янга, Р.В.\,Гамкрелидзе, Дж.\,Варги. Поэтому 


изучаемая в данной работе задача с поточечным фазовым ограничением


принципиально отличается от задач с аналогичным ограничением работ


\cite{sumin00b}, \cite{sumin00c}.





По этой причине


естественным представляется построение теории субоптимального


управления системами с распределенными параметрами, т.\,е. оптимизационной


теории для таких систем. Эта теория не основывается на факте существования


оптимальных элементов. В ней в качестве основного (искомого) объекта


рассматривается не оптимальный элемент, а {\it минимизирующая


последовательность} (м.\,п.) обычных допустимых управлений. Различные аспекты


теории субоптимального управления, в том числе и с поточечными фазовыми


ограничениями, в случае управляемых распределенных систем рассматривались ранее


в \cite{SuminD}, \cite{sumin00b}--\cite{sumin97b}.








В данной работе рассматривается так называемая параметрическая задача


субоптимального управления \cite{SuminD}, \cite{sumin00b}, \cite{sumin00c},


\cite{sumin97a}, \cite{sumin97b},


т.\,е. семейство зависящих от параметра задач субоптимального


управления.  Этот параметр является функциональным и аддитивно входит в


поточечное фазовое ограничение задачи. Изучение параметрических задач в


соответствии с общей идеологией метода возмущений (напр.,


\cite{altifom}, с.\,263) дает возможность получать информацию ``в целом''


о семействе и, как следствие, во многих важных частных случаях о каждой 


конкретной задаче отдельно.





В свою очередь, изучение параметрических задач субоптимального


управления естественным образом подводит к целесообразности


изучения в качестве м.\,п. так


называемых {\it минимизирующих приближенных решений} (м.\,п.\,р) в смысле


Дж.\,Варги \cite{warga}. Oни характеризуются, в частности, тем, что их элементы


удовлетворяют ограничениям лишь ``в пределе'' и


порождаемая именно таким понятием м.\,п. функция значений в самой общей ситуации


является полунепрерывной снизу\footnote{Заметим, что функция значений,


порождаемая классическим понятием м.\,п., элементы которой удовлетворяют


ограничениям в точном смысле, таким свойством, вообще говоря, не


обладает \cite{sumin97a}. Кроме того, важно отметить, что м.\,п.\,р. 


a)~всегда


существует, b)~позволяет записывать все результаты в терминах расширенной


\cite{gam77}, \cite{warga}


задачи, если такое расширение вообще возможно, c)~адекватно теории численных


методов, d)~удобно с прикладной (инженерной) точки зрения (подробности см.,


напр., в \cite{warga}).}.


%%%%%%%%%%%%%% 


Данное обстоятельство позволяет применить к


исследованию оптимизационных задач аппарат негладкого анализа,


а именно, анализа нормалей (проксимальных нормалей, нормалей Фреше) к замкнутым


множествам в банаховых пространствах и обобщенного дифференцирования негладких


(полунепрерывных снизу) функций в банаховых пространствах (напр.,


\cite{clarke}--\cite{mordukh}).





Библиография по теории оптимального управления системами


Гурса--Дарбу насчитывает десятки работ (напр.,


\cite{viplotnikovvisumin}, \cite{irkutyane}).


При этом задачи с поточечными фазовыми


ограничениями рассматривались, по-видимому, лишь в монографии \cite{irkutyane}


и только с точки зрения классических условий оптимальности.


Применяемый в данной работе подход \cite{SuminD}, \cite{sumin00b},


\cite{sumin00c}, \cite{sumin97a}, \cite{sumin97b}


позволяет изучить много


вопросов теории оптимизации, которые ранее не рассматривались в оптимальном


управлении гиперболическими уравнениями и, в частности, системами Гурса--Дарбу.


К ним относятся вопросы, связанные с необходимыми и достаточными


условиями для м.\,п.\,р., с необходимыми и достаточными условиями регулярности,


нормальности задач, с дифференциальными свойствами функций значений


(чувствительность) и др. Из-за ограниченности объема статьи здесь


рассматривается только часть этих вопросов.








\section{Постановка задачи}





Рассмотрим параметрическую задачу оптимального управления системой Гурса--Дарбу


с поточечным фазовым ограничением


\begin{equation}


I_0(u)\to\inf,\ \ I_1(u)\in


{\cal M}+q,\ \ u\in {\cal D},\quad q\in C(X)


\text{ --- параметр},\tag{$P_q$}


\end{equation}


где ${\cal D}\equiv\{u\in L_\infty^m(\Pi):u(x,\,y)\in


U$ п.\,в. в $\Pi\}$, $U\subset R^m$ --- компакт,


$\Pi\equiv[0,a]\times[0,b]$, $(\ov{x}_j, \ov{y}_j)$, $j=1,\dotsc,l$,


--- фиксированный набор точек, ${\cal M}$ --- множество всех непрерывных


неположительных на компакте $X\subseteq\Pi$ функций,


$I_0(u)\equiv\sum\limits_{j=1}^lG(z[u](\ov{x}_j, \ov{y}_j))$,


$I_1(u)(\cdot, \cdot)\equiv\Phi(\cdot, \cdot, z[u](\cdot, \cdot))$,


$z[u]$ --- соответствующее управлению $u\in{\cal D}$ абсолютно 


непрерывное решение нелинейного


векторного гиперболического уравнения ($z\equiv(z_1,\dotsc,z_n)\in R^n$,


$z_x\equiv p$,  $z_y\equiv r$)


\begin{equation}\label{equat}


z_{xy}=f(x,y,z,z_x,z_y,u(x,y)),\quad  z(x,0)=\varphi_1(x),\quad


z(0,y)=\varphi_2(y).


\end{equation}


Предположим, что функции


$f: \Pi\times R^n\times R^n\times R^n\times U\to R^n$, $G: R^n\to R$


вместе с якобианами $\nabla_zf$, $\nabla_pf$,


$\nabla_rf$ и градиентом $\nabla_zG$ измеримы в смысле Лебега по $(x,y)$


и непрерывны по $(z,p,r,v)$; функция $\Phi: \Pi\times R^n\to R$


и градиент $\nabla_z\Phi$ непрерывны по $(x,y,z)$; функции


$\varphi_1 :[0, a]\to R^n$ и $\varphi_2 :[0, b]\to R^n$ липшицевы, причем


$\varphi_1(0)=\varphi_2(0)$. Предположим также, что справедлива оценка


\begin{equation}


\label{eq12}


|f|+|\nabla_zf|+|\nabla_pf|+|\nabla_rf|\leq N(M)  


\ \forall


(x,y,z,p,r,v)\in\Pi\times S^n_M\times S^n_M\times S^n_M\times U,


\end{equation}


где $S^n_M\equiv\{x\in R^n : |x|<M\}$, $N(M)$ --- некоторая


функция аргумента $M>0$,


и для любого управления $u\in{\cal D}$


существует единственное абсолютно непрерывное


решение $z[u]$ задачи (\ref{equat}), причем


\begin{equation}\label{eval}


|z[u]|^{(0)}_\Pi+\|z_x[u]\|_{\infty, \Pi}+\|z_y[u]\|_{\infty, \Pi}\leq K  


\ \; \forall  u\in{\cal D},


\end{equation}


где $K>0$ --- постоянная, не зависящая от управления


$u\in{\cal D}$. Здесь и ниже используются следующие обозначения:


$\|\varphi\|_{p, \Pi}$ --- норма в пространстве $L^n_p(\Pi)$ суммируемых 


на $\Pi$ с $p$-й степенью (существенно ограниченных при $p=\infty$) функций


$\varphi :\Pi\to R^n$; $|\varphi|^{(0)}_\Pi$ --- норма в пространстве $C_n(\Pi)$


непрерывных на $\Pi$ функций $\varphi :\Pi\to R^n$.








\begin{notes}


Условия на правую часть $f$ уравнения (\ref{equat}), достаточные для


выполнимости требований \eqref{eq12}, \eqref{eval}, можно найти, например, в


\cite{plotsumv72}.


\end{notes}





Введем на множестве ${\cal D}$ метрику Экланда


$d(u^1, u^2)\equiv \mes\{(x, y)\in\Pi : u^1(x, y)\ne u^2(x, y)\}$, превратив его


тем самым в полное метрическое пространство. Согласно \cite{warga}


минимизирующим приближенным решением (м.\,п.\,р.) в задаче $(P_q)$ называется


такая последовательность управлений $u^i\in{\cal D}$, $i=1,2,\dotsc$, что


\begin{equation}\label{mas}


I_0(u^i)\leq\beta(q)+\gamma^i,\ \  u^i\in {\cal D}^{\varepsilon^i}_q,\ \  


\gamma^i,\varepsilon^i\geq 0,\ \  \gamma^i,\varepsilon^i\to 0,  


\ \ i\to\infty,


\end{equation}


где


${\cal D}^\varepsilon_q


\equiv\{u\in {\cal D} : \rho(I_1(u), {\cal M}+q)\leq\varepsilon\}$,


$\rho(I, {\cal M}+q)\equiv\inf\limits_{x\in {\cal M}}|x+q-I|^{(0)}_X$,


$\beta(q)\equiv\beta_{+0}(q)\equiv\lim\limits_{\varepsilon\to+0}


\beta_\varepsilon(q)\leq\beta_0(q)$, 


$\beta_\varepsilon(q)\equiv\{\inf\limits_{u\in {\cal D}^\varepsilon_q}I_0(u)$,


если


${\cal D}^\varepsilon_q\ne\emptyset$; 


$+\infty$, если ${\cal D}^\varepsilon_q=\emptyset\}$.


При этом функцию $\beta : C(X)\to R\cup\{+\infty\}$


принято называть функцией значений задачи $(P_q)$. Функция $\beta$ является


полунепрерывной снизу~\cite{sumin00c}.








\section{Принцип максимума}\label{outlook}





Введем обозначения:


$H(x,y,z,p,r,v,\eta)=\langle\eta, f(x,y,z,p,r,v)\rangle$,


$\xi\equiv(z, z_x, z_y)$,


\begin{gather*}


\Delta_uf(\xi[u]; u^1(\cdot,\cdot),u^2(\cdot,\cdot))\equiv


f(\cdot,\cdot,\xi[u^1](\cdot,\cdot), u^1(\cdot,\cdot))-


f(\cdot,\cdot,\xi[u^2](\cdot,\cdot),u^2(\cdot,\cdot)),\\


%


\nabla_zH[u](x, y, \eta)\equiv\nabla_zH


(x, y,z[u](x, y), z_x[u](x, y), z_y[u](x, y), u(x, y), \eta),


\end{gather*}


и т.\,п., $\langle\cdot,\cdot\rangle$ --- знак скалярного произведения. Обозначим


также через $\chi_{x, y}(\cdot,\cdot)$ характеристическую функцию


прямоугольника $[0, x]\times[0, y]$. 


Из результатов работы~\cite{plotsumv72} непосредственно вытекает








\begin{lems}\label{ev1}


Для любых двух управлений


$v,  u\in {\cal D}$ справедлива оценка


\begin{equation*} %%\label{ev2}


|z[v]-z[u]|^{(0)}_\Pi+\|z_x[v]-z_x[u]\|_{p, \Pi}+\|z_y[v]-z_y[u]\|_{p, \Pi}\leq


C_1\|\Delta_uf(\xi[u]; u^1(\cdot,\cdot), u^2(\cdot,\cdot))\|_{p, \Pi},


\end{equation*}


где $C_1>0$ зависит лишь от $p$, $1\leq p<\infty$.





Для любого управления $u\in {\cal D}$ существует единственное в $L^n_1(\Pi)$


решение~$\eta_0[u]$ сопряженной задачи


\begin{multline*} %%\label{adjoint}


\eta(x,y)=\int_x^a\int_y^b\nabla_z^*H[u](\xi_1,\xi_2,\eta(\xi_1,\xi_2))


d\xi_1 d\xi_2+


\int_y^b\nabla_p^*H[u](x,\xi_2,\eta(x,\xi_2)) d\xi_2+ \\


%


+\int_x^a\nabla_r^*H[u](\xi_1,y,\eta(\xi_1,y)) d\xi_1-


\sum_{j=1}^l\nabla_zG(z[u](\ov{x}_j,\ov{y}_j))\chi_{\ov{x}_j,\ov{y}_j}(x,y).


\end{multline*}


Это решение равномерно ограничено в $L_\infty^n(\Pi)$


и равномерно непрерывно в норме $L_p^n(\Pi)$ зависит от $u\in{\cal D}$ при любом


$1\leq p<\infty$. Сверх того, функционал $I_0: {\cal D}\to R$


и оператор $I_1: {\cal D}\to C(\Pi)$


также равномерно непрерывны.


\end{lems}





Пусть $M(\Pi)$ --- пространство мер Радона, сосредоточенных на


$\Pi$, а $\|\lambda\|$ --- полная вариация меры


$\lambda\in M(\Pi)$. 








\begin{lems}[{\series{m}\shape{n}\selectfont \cite{SuminD}}]


\label{adj2}


Для любой меры Радона $\lambda\in M(\Pi)$


и любого управления $u\in {\cal D}$


существует единственное в классе $L_1^n(\Pi)$ решение $\eta_1[u,\lambda]$


уравнения


\begin{multline}\label{adj22}


\eta(x,y)=\int_x^a\int_y^b\nabla_z^*H[u](\xi_1,\xi_2,\eta(\xi_1,\xi_2))


d\xi_1 d\xi_2+


\int_y^b\nabla_p^*H[u](x,\xi_2,\eta(x,\xi_2)) d\xi_2+ \\


%


+\int_y^b\nabla_r^*H[u](\xi_1,y,\eta(\xi_1,y)) d\xi_1-


\int_x^a\int_y^b\nabla_z\Phi(\xi_1,\xi_2, z[u](\xi_1,\xi_2)) 


\lambda(d\xi_1 d\xi_2),


\end{multline}


причем $\eta\in L_\infty^n(\Pi)$. 


Это решение равномерно непрерывно в норме $L_1^n(\Pi)$


зависит от $u\in {\cal D}$, $\lambda\in M(\Pi)$ и справедлива оценка


$\|\eta_1[u, \lambda]\|_{\infty, \Pi}\leq C_2\|\lambda\|$,


в которой $C_2>0$ не зависит от $u\in{\cal D}$, $\lambda\in M(\Pi)$.


\end{lems}








\begin{thms}[{\series{m}\shape{n}\selectfont принцип максимума для м.\,п.\,р.}]


\label{Pontryagin}


Пусть $\beta(q)<+\infty$. Если $u^k\in{\cal D}$, $k=1,2,...$,  %%% return \dotsc


есть м.\,п.\,р.


в задаче $(P_q)$, то найдутся сходящаяся к нулю последовательность


чисел $\gamma^k\geq 0$, последовательность чисел $\mu^k_0\geq 0$ и


последовательность сосредоточенных на множестве


$\Pi_k\equiv\{(x,y)\in X : |\Phi(x,y,z[u^k](x,y))-q(x,y)|\leq\gamma^k\}$


неотрицательных атомических мер Радона $\lambda^k\in M(\Pi)$,


$k=1,2,\dotsc$, такие, что


\begin{gather}\label{normir}


u^k\in {\cal D}^{\gamma^k}_q,\quad  \mu^k_0+\|\lambda^k\|=1, \\


%


\int_\Pi\{\max_{v\in U}H(x,y,\xi[u^k](x,y), v, \eta[u^k, \lambda^k, 


\mu_0^k](x,y))-\notag \\


%


\label{maxcond}


-H(x,y,\xi[u^k](x,y),u^k(x,y),\eta[u^k,


\lambda^k, \mu_0^k](x,y) )\} dx\,dy\leq\gamma^k,


\end{gather}


где $\eta[u, \lambda, \mu]\equiv\mu\eta_0[u]+\eta_1[u, \lambda]$.


\end{thms}








\begin{notes}


Из теоремы~\ref{Pontryagin} следует, что если управление $u_0\in {\cal D}^0_q$


таково, что $I_0(u)=\beta(q)$, то оно удовлетворяет обычному принципу максимума


($\mu^k_0\equiv\mu_0$, $\lambda^k\equiv\lambda$, $\gamma^k\equiv\gamma=0$). 


Несколько видоизменяя рассуждения приводимого ниже доказательства 


теоремы~\ref{Pontryagin}, можно показать, что такой же обычный принцип 


максимума справедлив и для такого


управления $u_0\in {\cal D}^0_q$, что $I_0(u)=\beta_0(q)$.


\end{notes}








Приведем кратко схему доказательства принципа максимума


\cite{sumin1986}. Пусть $u^k$, $k=1,2,\dotsc$, --- м.\,п.\,р.


в задаче $(P_q)$. Рассмотрим вспомогательную задачу


\begin{equation}\label{helpproblem}


J(u)\to\inf,\quad  u\in{\cal D},


\end{equation}


где $J(u)\equiv\max\{I_0(u)-\beta(q)$,  


$\Phi(x, y, z[u](x, y))-q(x, y)$, $(x, y)\in X\}$.


Очевидно, последовательность $u^k$ является


одновременно минимизирующей и для задачи (\ref{helpproblem}) с нулевой


нижней гранью. Пусть


$X_k\equiv\{(x_{k, j}, y_{k, j}) : j=1,\dotsc,l_k\}\subset X$ ---


$1/k$-сеть в $X$. Рассмотрим семейство вспомогательных задач


$J_k(u)\to\inf$,  $u\in {\cal D}$, где


$J_k(u)\equiv\max\{I_0(u)-\beta(q)$;  $\Phi(x, y, z[u](x, y))-q(x, y)$,  


$(x, y)\in X_k\}$. В силу леммы \ref{ev1}


функционал $J_k(\cdot)$ непрерывен и ограничен на ${\cal D}$ и,


как нетрудно показать, пользуясь равномерной ограниченностью и


равностепенной непрерывностью семейства $\{z[u] : u\in {\cal D}\}\subset


C(\Pi)$, имеет место предельное соотношение


$$


\lim\limits_{k\to\infty}\inf\limits_{u\in {\cal D}}J_k(u)= \inf\limits_{u\in


{\cal D}}J(u)=\lim\limits_{k\to\infty}J(u^k)=0.


$$





Отсюда следует, что найдется такая последовательность ${\delta^k\geq 0}$,


$k=1,2,\dotsc$, $\delta^k\to 0$, $k\to\infty$, что


$J_k(u^k)\leq\inf\limits_{u\in {\cal D}}J_k(u)+\delta^k$. Поэтому в силу


вариационного принципа Экланда \cite{ekeland} существует


последовательность $v^k\in{\cal D}$, элементы которой являются 


решениями соответствующих задач


\begin{equation}\label{jkproblem}


J_k(u)+\sqrt{\delta^k}d(v^k, u)\to\min,\quad  u\in {\cal D},


\end{equation}


причем справедливы неравенства


\begin{equation}\label{closeness}


d(u^k, v^k)\leq\sqrt{\delta^k},\quad  J_k(v^k)\leq J_k(u^k).


\end{equation}


В отличие от $J(u)$ функционал~$J_k(u)$


представляет собой максимум лишь от конечного числа


функционалов. Поэтому необходимые условия оптимальности управления~$v^k$ в


задаче (\ref{jkproblem}) устанавливаются подобно тому, как это обычно делается в


задачах оптимального управления с конечным числом функциональных ограничений


(напр., \cite{sumqual}). Последнее обстоятельство позволяет 


говорить о существовании такого вектора


$\mu^k\equiv(\mu^k_0, \mu^k_1,\dotsc,\mu^k_{l_k})\in R^{l_k+1}$, что


\begin{gather}\label{mu1}


\mu^k_j\geq 0,\quad  j=0,\dotsc,l_k;\quad  \sum_{j=0}^{l_k}\mu^k_j=1,\\


%


\label{mu2}


\mu^k_j(J_k(v^k)-(\Phi(x_{k, j}, y_{k, j}, z[v^k](x_{k, j}, y_{k, j}))-


q(x_{k, j}, y_{k, j}))),\quad j=1,\dotsc,l_k,\\


%


H\bigg(x, y, \xi[v^k](x, y), v^k(x, y),\mu^k_0\eta_0[v^k](x, y)+


\sum_{j=1}^{l_k}\mu^k_j\wh{\eta}[v^k](x, y, x_{k, j}, y_{k, j})\bigg)-\notag\\


%


\label{mu3}


-H\bigg(x, y, \xi[v^k](x, y), w,\mu^k_0\eta_0[v^k](x, y)+


\sum_{j=1}^{l_k}\mu^k_j\wh{\eta}[v^k](x, y, x_{k, j}, y_{k, j})\bigg)


\geq-2\sqrt{\delta^k} \\


%


\forall w\in U\ \ \text{при п.\,в.}\ \  (x, y)\in\Pi,\notag


\end{gather}


где $\wh{\eta}[u](\cdot,\cdot,\ov{x},\ov{y})$ --- решение уравнения


\begin{multline}\label{adj3}


\eta(x,y,\ov{x},\ov{y})=


\int_x^{a}\int_y^{b}\nabla_z^*H[u](\xi_1,\xi_2,


\eta(\xi_1, \xi_2, \ov{x}, \ov{y})) d\xi_1 d\xi_2+ 


\int_y^{b}\nabla_p^*H[u](x, \xi_2, 


\eta(x, \xi_2, \ov{x}, \ov{y})) d\xi_2+ \\


%


+\int_x^{a}\nabla_r^*H[u](\xi_1, y, 


\eta(\xi_1, y, \ov{x}, \ov{y})) d\xi_1-


\nabla_z\Phi(\ov{x}, \ov{y},z[u](\ov{x}, \ov{y}))\chi_{\ov{x},\ov{y}}(x,y).


\end{multline}





Введем посредством равенства $\lambda^k\equiv$


$\sum\limits_{j=1}^{l_k}\mu_j^k\delta_{(x_{k,j},y_{k,j})}$


атомические меры Радона $\lambda^k\in M(\Pi)$,


где $\delta_{(x, y)}$ --- $\delta$-мера Дирака,


сосредоточенная в точке $(x, y)$.





Тогда $\sum\limits_{j=1}^{l_k}\mu_j^k\wh{\eta}[v^k](x, y, x_{k, j}, y_{k, j})=


\int\limits_\Pi\wh{\eta}[v^k](x, y, \ov{x}, \ov{y})\lambda^k(d\ov{x}\,d\ov{y})$.


Обозначая последнюю функцию через $\Theta[v^k,\lambda^k]$ и интегрируя


уравнение~(\ref{adj3}) по мере $\lambda^k$, выводим


\begin{multline*}


\Theta[v^k,\lambda^k](x, y)=


\displaystyle\int_x^{a}\int_y^{b}\nabla_z^*H[u](\xi_1, \xi_2, 


\Theta[v^k,\lambda^k](\xi_1, \xi_2)) d\xi_1 d\xi_2+\\


%


+\int_y^{b}\nabla_p^*H[u](x, \xi_2, 


\Theta[v^k,\lambda^k](x, \xi_2)) d\xi_2+ \\


%


+\int_x^{a}\nabla_r^*H[u](\xi_1, y, 


\Theta[v^k,\lambda^k](\xi_1, y)) d\xi_1-


\int_x^a\int_y^b\nabla_z\Phi(\ov{x},


\ov{y},z[u](\ov{x}, \ov{y}))\lambda^k(d\ov{x}\,d\ov{y}),


\end{multline*}


т.\,е. $\Theta[v^k,\lambda^k]$ --- решение уравнения (\ref{adj22}) при $u=v^k$,


$\lambda=\lambda^k$, которое согласно лемме \ref{adj2} единственно в


классе $L_1^n(\Pi)$. Следовательно,


$\Theta[v^k,\lambda^k](x,y)\equiv\eta_1[v^k,\lambda^k](x,y)$, и


соотношение (\ref{mu3}) может быть переписано в виде


\begin{gather}


H(x,y,\xi[v^k](x,y),v^k(x,y),\eta[v^k,\lambda^k,\mu^k_0](x,y))- \notag \\


%


\label{mu03}


-H(x,y,\xi[v^k](x,y),w,\eta[v^k,\lambda^k,\mu^k_0](x,y))


\geq-2\sqrt{\delta^k}\ \;  \forall w\in U\ \ \text{при п.\,в.}


\ \ (x, y)\in\Pi.


\end{gather}


С помощью условия близости (\ref{closeness}) двух управлений $u^k$


и $v^k$, а также априорных оценок лемм \ref{ev1}, \ref{adj2} соотношение


(\ref{mu03}) может быть переписано в терминах исходного м.\,п.\,р.


$u^k$, $k=1,2,\dotsc$ Опуская весьма громоздкие,


но достаточно очевидные детали такого переписывания,


получаем соотношение максимума (\ref{maxcond}). Равенство же (\ref{mu1})


переходит в условие невырожденности набора множителей


Лагранжа (\ref{normir}). Используя непрерывность на $X$ функции $q$, близость


управлений $u^k$ и $v^k$ и равномерную непрерывность по $u\in {\cal D}$ решений


$z[u]$, из (\ref{mu2}) заключаем, что $\lambda^k$


сосредоточена на множестве $\Pi_k$.$\quad\square$


\medskip





В заключение этого раздела, как и в \cite{sumin00b}, введем 








\begin{defns}


Последовательность $u^k\in{\cal D}$, $k=1,2,\dotsc$, назовем стационарной в


задаче $(P_q)$, если найдутся последовательность неотрицательных чисел


$\gamma^k\to0$, $k\to\infty$, $u^k\in {\cal D}^{\gamma^k}_q$, и


ограниченная последовательность пар $(\mu^k_0,\lambda^k)$, $\mu_0^k\geq0$,


$\lambda^k\in M(\Pi)$ --- неотрицательная атомическая мера Радона,


сосредоточенная на множестве $\Pi_k$, такие, что выполняется


неравенство (\ref{maxcond}) и все предельные точки последовательности


пар $(\mu^k_0,\lambda^k)$, $k=1,2,\dotsc$ (в $*$-слабом смысле для второй


компоненты), отличны от нуля. Стационарную в задаче $(P_q)$ последовательность


$u^k\in{\cal D}$, $k=1,2,\dotsc$, назовем нормальной (регулярной,


анормальной), если все (существуют, не существуют) предельные точки любой


(хотя бы одной, ни у одной) соответствующей ей последовательности


$\mu^k_0$ не равны нулю. Задача $(P_q)$ называется


нормальной (регулярной, анормальной), если все (существуют, не


существуют) ее стационарные последовательности нормальны (являющиеся


регулярными, являющиеся регулярными).


\end{defns}








\section{Аппроксимация задачи с фазовыми ограничениями,\protect\\


липшицевость функции значений}





Рассмотрим последовательность семейств


оптимизационных задач, зависящих от конечномерного векторного параметра


$q^k\equiv(q^k_1,\dotsc,q^k_{l_k})\in R^{l_k}$, аппроксимирующих исходное


семейство $(P_q)$:


\begin{equation}


I_0(u)\to\inf,\ \  I^k(u)\in {\cal


M}^k+q^k,\ \  u\in{\cal D},\quad


q^k\in R^{l_k} \text{ --- параметр},\tag{$P^k_{q^k}$}


\end{equation}


где ${\cal M}^k\equiv\{y\in


R^{l_k}:y_i\leq 0$,  $i=1,\dotsc,l_k\}$,


$I^k(u)\equiv(I^k_1(u),\dotsc,I^k_{l_k}(u))$,


$$


I^k_i(u)\equiv\Phi(x_{k, i},y_{k, i},z[u](x_{k, i},y_{k, i})),\ \


(x_{k, i}, y_{k, i})\in X_k,


$$


а $1/k$-сеть $X_k$ определена в


разделе~\ref{outlook}. Как и в случае задачи $(P_q)$, согласно \cite{warga}


м.\,п.\,р. в задаче $(P^k_{q^k})$ называется такая последовательность управлений


$u^i\in {\cal D}$, $i=1,2,\dotsc$, что выполняются соотношения


\begin{equation*}


I_0(u^i)\leq\beta_k(q^k)+\gamma^i,\ \ u^i\in {\cal D}^{k,\varepsilon^i}_{q^k},  


\ \ \gamma^i,\varepsilon^i\geq 0,\ \ \gamma^i, \varepsilon^i\to0,\ \ i\to\infty,


\end{equation*}


где ${\cal D}^{k, \varepsilon}_{q^k}\equiv\{u\in {\cal D}:


I^k_j(u)\leq q^k_j+\varepsilon$, $j=1,\dotsc,l_k\}$,


$\beta_k(q^k)\equiv\beta_{k, +0}(q^k)\equiv


\lim\limits_{\varepsilon\to+0}\beta_{k, \varepsilon}(q^k)\leq\beta_{k, 0}(q^k)$;


$\beta_{k, \varepsilon}(q^k)\equiv\{\inf


\limits_{u\in {\cal D}^{k, \varepsilon}_{q^k}}I_0(u)$,


${\cal D}^{k, \varepsilon}_{q^k}\neq\emptyset$; $+\infty$,


${\cal D}^{k, \varepsilon}_{q^k}=\emptyset\}$. Можно утверждать, что функция


значений $\beta_k : R^{l_k}\to R$, как и функция $\beta$, полунепрерывна снизу


\cite{sumin00c}.








\begin{lems}\label{app}


Пусть $\beta(q)<+\infty$, $q\in C(X)$. Тогда существует такая


последовательность векторов $q^k\in R^{l_k}$, $k=1,2,\dotsc$, что


$\beta_k(q^k)\to\beta(q)$, $k\to\infty$. В частности, таковой


является последовательность


$\ov{q}^k\equiv(\ov{q}^k_1,\dotsc,\ov{q}^k_{l_k})$,


$\ov{q}^k_i=q(x_{k, i},y_{k, i})$, $i=1,\dotsc,l_k$.


\end{lems}





Доказательство этой аппроксимационной леммы аналогично 


доказательству леммы 3.1 из \cite{sumin00c} и поэтому опускается. Следует


отметить, что данная лемма может быть полезна при конструировании


численных методов решения задач с фазовыми ограничениями.





Введем так называемые обычный и сингулярный обобщенные субдифференциалы в


смысле \cite{mordukhovich88} для произвольной полунепрерывной снизу функции


$\Theta : R^n\to R\cup\{+\infty\}$ в каждой точке $q\in \dom\Theta$


соответственно по формулам $\partial\Theta(q)\equiv


D^-\Theta(q)\equiv\{\xi : (\xi, -1)\in K((q,\Theta(q))\mid\epi\Theta)\}$,


$\partial^\infty\Theta(q)\equiv D^-_s\Theta(q)\equiv\{\xi : (\xi, 0)\in


K((q,\Theta(q))\mid\epi\Theta)\}$, где


$$


K((q,\Theta(q))\mid\epi\Theta)\equiv


\limsup_{(r,\sigma)\to(q,\Theta(q))}P((r,\sigma)\mid\epi\Theta)


$$


называется конусом обобщенных нормалей к


$\epi\Theta$ в точке $(q,\Theta(q))$,


\begin{align*}


P((r,\sigma)\mid\epi\Theta)&\equiv


\{\alpha((r,\sigma)-(s,\nu)):\alpha>0,


\ (s,\nu)\in W((r,\sigma)\mid\epi\Theta)\},\\


W((r,\sigma)\mid\epi\Theta)&\equiv\{(s,\nu)\in \epi\Theta: 


|(r,\sigma)-(s,\nu)|=\rho((r,\sigma)\mid\epi\Theta)\},


\end{align*}


$\rho((r,\sigma)\mid\epi\Theta)$ --- евклидово расстояние от точки $(r,\sigma)$


до множества $\epi\Theta$ (\cite{mordukhovich88}, с.\,45).


Покажем, что справедлива следующая лемма о связи множителей Лагранжа с


обобщенными нормалями в смысле \cite{mordukhovich88}.








\begin{lems}\label{lstat}


Пусть $\beta_k(q^k)<\infty$ и $(\zeta,-\tau)\in


K((q^k,\beta_k(q^k))\mid\epi\beta_k)$ --- произвольная обобщенная нормаль,


причем $(\zeta,-\tau)\neq 0$.


Тогда существуют последовательность неотрицательных чисел $\gamma^i\to 0$,


$i\to\infty$, ограниченная последовательность множителей Лагранжа


$\mu^i\equiv(\mu^i_0,\dotsc,\mu^i_{l_k})\in R^{l_k}$ и последовательность


управлений $u^i\in {\cal D}^{k,\gamma^i}_{q^k}$, $i=1,2,\dotsc$, такие, что


\begin{gather}\label{muic}


|\mu^i|\neq 0,\ \ \mu^i_j\geq 0,\ \ j=0,\dotsc,l_k;\ \  \mu^i_j=0,


\text{ если } |I^k_j(u^i)-q^k_j|\geq\gamma^i,\ \ j=1,\dotsc,l_k, \\


%


\int_\Pi\{\max_{v\in U}H(x,y,\xi[u^i](x,y),v,\eta[u^i,\lambda^i,\mu^i_0]


(x,y))-\notag\\


%


\label{maxc}


-H(x,y,\xi[u^i](x,y),u^i(x,y),\eta[u^i,\lambda^i,\mu^i_0](x,y))\} dx\,dy


\leq\gamma^i; \\


%


%% \label{kmi}


\zeta+\sum_{j=1}^{l_k}\mu_je^j=0, \notag


\end{gather}


где $\lambda^i\equiv\sum\limits_{j=1}^{l_k}\mu^i_j\delta_{(x_{k,j},y_{k,j})}$,


а $\mu\equiv(\tau,\mu_1,\dotsc,\mu_{l_k})\ne 0$ --- предельная


точка последовательности векторов $\mu^i$, $i=1,2,\dotsc$


\end{lems}








{\bf Доказательство.}


Пусть $(r^{k,i},\sigma^{k,i})\notin \epi\beta_k$, $i=1,2,\dotsc$, ---


произвольная последовательность точек, сходящаяся к 


$(q^k,\beta_k(q^k))$. Тогда из полунепрерывности снизу функции $\beta_k$ и


определения конуса $K((q^k,\beta_k(q^k))\mid\epi\beta_k)$


(\cite{mordukhovich88}, с.\,45) следует существование последовательности


положительных чисел $\alpha^i$ и последовательности точек


\begin{equation}\label{*}


(q^{k,i},\beta_k(q^{k,i}))\in W((r^{k,i},\sigma^{k,i})\mid\epi\beta_k)


\subset R^{l_k+1},\quad  i=1,2,\dotsc,


\end{equation}


таких, что


\begin{equation}\label{**}


q^{k,i}\to q^k,\ \ \beta_k(q^{k,i})\to\beta_k(q^k),  


\ \ \alpha^i(\zeta^{k,i},-\tau^{k,i})\to(\zeta,-\tau),\ \  i\to\infty,


\end{equation}


$(\zeta^{k,i},-\tau^{k,i})\equiv


(r^{k,i},\sigma^{k,i})-(q^{k,i},\beta_k(q^{k,i}))$, $\tau^{k,i}\geq 0$.


В множестве $W((r^{k,i},\sigma^{k,i})\mid\epi\beta_k)$ точка


$(q^{k,i},\beta_k(q^{k,i}))$ выбирается произвольно.


Для дальнейшего доказательства леммы покажем, что справедлива








\begin{lems}\label{amsa} Пусть $u^{i,s}$, $s=1,2,\dotsc$, ---


произвольное м.\,п.\,р. в смысле $(\ref{mas})$ в задаче $(P^k_{q^{k,i}})$.


Тогда существуют последовательность чисел $\gamma^{i,s}\geq 0$,


$\gamma^{i,s}\to 0$, $s\to\infty$,


$u^{i,s}\in {\cal D}^{k, \gamma^{i,s}}_{q^{k,i}}


\equiv\{u\in{\cal D} : I^k_j(u)\leq q^{k,i}_j+\gamma^{i,s}$,


$j=1,\dotsc,l_k\}$,


последовательность векторов


$\mu^{i,s}\equiv(\mu_0^{i,s},\dotsc,\mu_{l_k}^{i,s})\in R^{l_k+1}$,


$s=1,2,\dotsc$,


\begin{equation}\label{***}


\sum_{j=0}^{l_k}\mu^{i, s}_j=1,\ \  


\mu^{i,s}_j\geq 0,\ \ j=0,\dotsc,l_k; 


\ \ \mu^{i,s}_j=0 \text{ при } |I^k_j(u^{i,s})-q^{k,i}_j|\geq\gamma^{i,s},


\ \ j=1,\dotsc,l_k,


\end{equation}


такие, что


\begin{gather}


\int_{\Pi}\{\max_{v\in U}H(x,y,\xi[u^{i,s}](x,y),v,


\eta[u^{i,s}, \lambda^{i,s}, \mu^{i,s}_0\tau^{k,i}](x,y))-\notag \\


%


\label{****}


-H(x,y,\xi[u^{i,s}](x,y),u^{i,s}(x,y),\eta[u^{i,s},\lambda^{i,s},


\mu^{i,s}_0\tau^{k,i}](x,y))\} dx\, dy\leq\gamma^{i,s}, \\


%


\label{*****}


\max_{q\in\ov{S}^{l_k}_P}\Big\langle-


\mu_0^{i,s}\zeta^{k,i}-\sum_{j=1}^{l_k}\mu_j^{i,s}e^j,


q-q^{k,i}\Big\rangle\leq\gamma^{i,s},


\end{gather}


где


$\lambda^{i, s}\equiv\sum\limits_{j=1}^{l_k}\mu^{i, s}_j


\delta_{(x^{k, j}, y^{k, j})}$,


$e^j\equiv(\,\underbrace{0,\dotsc,0}_{j-1},1,0,\dotsc,0)\in R^{l_k}$,


а $P>0$ --- произвольное достаточно большое число, например,


$P=\sup\limits_{i\geq1}|q^{k,i}|+1$.


\end{lems}








\begin{pf}


В силу включения (\ref{*}) вектор $(\zeta^{k,i},-\tau^{k,i})$ является


перпендикуляром (\cite{clarke}, с.\,66) к $\epi\beta_k$ в точке


$(q^{k,i},\beta_k(q^{k,i}))$. Поэтому на основании предложения 2.5.5 из


\cite{clarke} выводим


$$


\langle(\zeta^{k,i},-\tau^{k,i}),(q^\prime,I)-


(q^{k,i},\beta_k(q^{k,i}))\rangle\leq


\frac{1}{2}|(q^\prime,I)-(q^{k,i},\beta_k(q^{k,i}))|^2  


\ \ \forall q^\prime\in R^{l_k},\ \; I\geq\beta_k(q^\prime),


$$


откуда после элементарных преобразований получаем


$$


\tau^{k,i}\beta_k(q^{k,i})-\langle\zeta^{k,i},q^{k,i}\rangle


\leq\tau^{k,i}I-


\langle\zeta^{k,i},q^\prime\rangle+\frac{1}{2}|q^\prime-q^{k,i}|^2+


\frac{1}{2}(I-\beta_k(q^{k,i}))^2  


\ \ \forall q^\prime\in R^{l_k},\ \; I\geq\beta_k(q^\prime).


$$


Докажем, что если последовательность $u^{i,s}$, $s=1,2,\dotsc$, ---


м.\,п.\,р. в задаче $(P^k_{q^{k,i}})$, то


последовательность пар $(u^{i,s},q^{k,i})$, $s=1,2,\dotsc$, ---


м.\,п.\,р. в задаче


\begin{equation}\label{pairprob}


I^i(u,q^\prime)\to\inf,\ \  I^k(u)-q^{\prime}\in {\cal M}^k,  


\ \ (u,q^\prime)\in\wh{{\cal D}}\equiv{\cal D}\times\ov{S}_P^{l_k},


\end{equation}


где


$I^i(u,q^\prime)\equiv\tau^{k,i}I_0(u)-\langle\zeta^{k,i},q^\prime\rangle+


\frac{1}{2}|q^\prime-q^{k,i}|^2+\frac{1}{2}(I_0(u)-\beta_k(q^{k,i}))^2$.


Прежде всего, наделим $\wh{{\cal D}}$ метрикой


$\wh{d}((u^1,q^1),(u^2,q^2))\equiv d(u^1,u^2)+|q^1-q^2|$ и введем следующие


обозначения: $\wh{{\cal D}}^t\equiv


\{(u,q^\prime)\in\wh{{\cal D}} : I^k_j(u)-q^\prime_j


\leq t$,  $j=1,\dotsc,l_k\}$, $\wh{\beta}_k^i\equiv


\lim\limits_{t\to0}\wh{\beta}_k^{i,t}$,


$\wh{\beta}_k^{i,t}\equiv\inf\limits_{\wh{{\cal D}}^t}I^i(u,q^\prime)$.


В соответствии с общим определением~\cite{warga} последовательность пар


$(\ov{u}^s,q^{\prime s})\in\wh{{\cal D}}$, $s=1,2,\dotsc$, будет м.\,п.\,р. в


задаче (\ref{pairprob}), если найдутся такие последовательности чисел


$\ov{\gamma}^s$, $\ov{\varepsilon}^s\geq0$, $\ov{\gamma}^s$,


$\ov{\varepsilon}^s\to0$, $s\to\infty$,


что


\begin{equation*} %%\label{mpao}


I^i(\ov{u}^s,q^{\prime s})\leq\wh{\beta}_k^i+\ov{\gamma}^s,  


\ \ (\ov{u}^s,q^{\prime s})\in\wh{{\cal D}}^{\ov{\varepsilon}^s},


\ \  s=1,2,\dotsc


\end{equation*}


Так как $u^{i,s}$, $s=1,2,\dotsc$, --- м.\,п.\,р. в задаче $(P^k_{q^{k,i}})$,


то найдутся такие последовательности неотрицательных чисел


$\varepsilon^{i,s}$, $\gamma^{i,s}\to0$,


$s\to\infty$, что $u^{i,s}\in {\cal D}^{k,\varepsilon^{i,s}}_{q^{k,i}}$,


$I_0(u^{i,s})\leq\beta_k(q^{k,i})+\gamma^{i,s}$. Поэтому


\begin{multline}\label{opp}


I^i(u^{i,s},q^{k,i})=\tau^{k,i}I_0(u^{i,s})-\langle\zeta^{k,i},q^{k,i}\rangle+


\frac{1}{2}(I_0(u^{i,s})-\beta_k(q^{k,i}))^2\leq\\


%


\le\tau^{k,i}\beta_k(q^{k,i})+\tau^{k,i}\gamma^{i,s}-


\langle\zeta^{k,i},q^{k,i}\rangle+\frac{1}{2}(\gamma^{i,s})^2.


\end{multline}


Поскольку нижняя грань $\wh{\beta}_k^i$ в задаче (\ref{pairprob}) равна


$\wh{\beta}_k^i=\tau^{k,i}\beta_k(q^{k,i})-\langle\zeta^{k,i},q^{k,i}\rangle$


(доказательство данного факта аналогично доказательству равенства (5.6) в


работе \cite{sumin97a}), то оценку (\ref{opp}) можно переписать в виде


\begin{equation}\label{opp1}


I^i(u^{i,s},q^{k,i})\leq\wh{\beta}_k^i+\wt{\gamma}^{i,s},


\end{equation}


где


$\wt{\gamma}^{i,s}\equiv\tau^{k,i}\gamma^{i,s}+\frac{1}{2}(\gamma^{i,s})^2$.


Кроме того, из соотношений $\varepsilon^{i,s}\to0$, $s\to\infty$,


$u^{i,s}\in {\cal D}^{k,\varepsilon^{i,s}}_{q^{k,i}}$ и определения множеств


$\wh{{\cal D}}^t$ следует, что


$(u^{i,s},q^{k,i})\in\wh{{\cal D}}^{\varepsilon^{i,s}}$. Последнее включение


в совокупности с оценкой~(\ref{opp1}) доказывает, что последовательность


пар $(u^{i,s},q^{k,i})$ действительно является м.\,п.\,р.


в задаче~(\ref{pairprob}), а для последовательности


$r^{i,s}\equiv\wh{\beta}_k^i-\wh{\beta}_k^{i,\varepsilon^{i,s}}+


\wt{\gamma}^{i,s}$, $i=1,2,\dotsc$, справедливы соотношения


$$


I^{i}(u^{i,s},q^{i,k})\leq\wh{\beta}_k^{i,\varepsilon^{i,s}}+r^{i,s},  


\quad (u^{i,s},q^{k,i})\in\wh{{\cal D}}^{\varepsilon^{i,s}}.


$$


Отсюда с помощью вариационного принципа Экланда \cite{ekeland} выводим


существование последовательности $(u^{i,s,1},q^{k,i,s})\in\wh{{\cal


D}}^{\varepsilon^{i,s}}$, являющейся решением в задаче


\begin{equation}\label{hm}


I^i(u,q^\prime)+\sqrt{r^{i,s}}\,\wh{d}((u,q^\prime),(u^{i,s,1},q^{k,i,s}))


\to\inf,  \quad (u,q^\prime)\in\wh{{\cal D}}^{\varepsilon^{i,s}},


\end{equation}


и такой, что


\begin{equation}\label{ea}


\wh{d}((u^{i,s},q^{k,i}),(u^{i,s,1},q^{k,i,s}))\leq\sqrt{r^{i,s}},  \quad


I^i(u^{i,s,1},q^{k,i,s})\leq I^i(u^{i,s},q^{k,i}).


\end{equation}


При этом из первой оценки (\ref{ea}) и леммы \ref{ev1} вытекает, что


$I_0(u^{i,s,1})-\beta_k(q^{k,i})\to0$, $s\to\infty$, и, как следствие,


последовательность $u^{i,s,1}$, $s=1,2,\dotsc$, является м.\,п.\,р. и в


задаче $(P^k_{q^k})$ при $q^k=q^{k,i}$. Из (\ref{hm}) можно заключить, что


\begin{equation*}


J_\gamma(u^{i,s,1},q^{k,i,s})\leq


\inf_{\wh{{\cal D}}}J_\gamma(u,q^\prime)+\gamma,


\end{equation*}


где


$J_\gamma(u,q^\prime)\equiv\max\{I^i(u,q^\prime)-I^i(u^{i,s,1},q^{i,k,s})


+\gamma+\sqrt{r^{i,s}}\wh{d}((u,q^\prime),(u^{i,s,1},q^{k,i,s}))$,


$I^k_j(u)-q^\prime_j-\varepsilon^{i,s}$, $j=1,\dotsc,l_k\}$.


Снова применяя вариационный принцип


Экланда, выводим, что найдется пара $(u^{i,s,1,\gamma},q^{k,i,s,\gamma})\in


\wh{{\cal D}}$,


являющаяся решением в задаче


\begin{equation}\label{ppp}


J_\gamma(u,q^\prime)+\sqrt{\gamma}\wh{d}((u,q^\prime),


(u^{i,s,1,\gamma},q^{k,i,s,\gamma}))\to\inf,


\quad (u,q^\prime)\in\wh{{\cal D}},


\end{equation}


и удовлетворяющая неравенствам


\begin{equation}\label{eaa}


\wh{d}((u^{i,s,1},q^{k,i,s}),(u^{i,s,1,\gamma},q^{k,i,s,\gamma}))\leq


\sqrt{\gamma},  \quad


J_\gamma(u^{i,s,1,\gamma},q^{k,i,s,\gamma})\leq J_\gamma(u^{i,s,1},q^{k,i,s}).


\end{equation}


Для получения необходимых условий оптимальности в последней задаче можно


воспользоваться результатами работы \cite{sumqual}, в которой достаточно


подробно изложен метод получения условий оптимальности в аналогичной задаче


для квазилинейного эллиптического уравнения (см. задачу для функционала


$J_{\xi,\rho}$ в \cite{sumqual}, с.\,1414). Применяя здесь этот метод,


получим следующее заключение. Так как


$(u^{i,s,1,\gamma},q^{k,i,s,\gamma})$ --- решение задачи (\ref{ppp}), то


существуют вектор


$\mu^{i,s,1,\gamma}\equiv(\mu^{i,s,1,\gamma}_0,\mu^{i,s,1,\gamma}_1,\dotsc,


\mu^{i,s,1,\gamma}_{l_k})\in R^{l_k+1}$ и достаточно большое число $K>0$


такие, что


\begin{gather*}


\mu^{i,s,1,\gamma}_j\geq 0,\ \  j=0,\dotsc,l_k,  


\ \ \sum\limits_{j=0}^{l_k}\mu^{i,s,1,\gamma}_j=1,\\


%


\mu^{i,s,1,\gamma}_j(J_\gamma(u^{i,s,1,\gamma},q^{k,i,s,\gamma})-


(I^k_j(u^{i,s,1,\gamma})-


q^{k,i,s,\gamma}-\varepsilon^{i,s}) )=0,\ \  j=1,\dotsc,l_k,\\


%


H(x,y,\xi[u^{i,s,1,\gamma}](x,y),u^{i,s,1,\gamma}(x,y),


\eta[u^{i,s,1,\gamma},\lambda^{i,s,1,\gamma},


\mu_0^{i,s,1,\gamma}\tau^{k,i}](x,y))-\\


%


-H(x,y,\xi[u^{i,s,1,\gamma}](x,y),w,\eta[u^{i,s,1,\gamma},


\lambda^{i,s,1,\gamma},\mu_0^{i,s,1,\gamma}\tau^{k,i}](x,y))\geq\\


%


\ge -2K(\sqrt{r^{i,s}}+\sqrt{\gamma}) 


\ \; \forall w\in U\text{ при п.\,в. } (x,y)\in\Pi,\\


%


\Big\langle\mu_0^{i,s,1,\gamma}(-\zeta^{k,i}+q^{k,i,s,\gamma}-q^{k,i})-


\sum\limits_{j=1}^{l_k}\mu^{i,s,1,\gamma}_je^j,q^\prime-q^{k,i,s,\gamma}


\Big\rangle


\leq 2K(\sqrt{r^{i,s}}+\sqrt\gamma)\ \; \forall q^\prime\in \ov{S}^{l_k}_P.


\end{gather*}


Переходя в полученном семействе необходимых условий к пределу при $\gamma\to


0$ и учитывая оценки (\ref{ea}), (\ref{eaa}), после весьма громоздких


вычислений приходим к выводу о справедливости 


леммы~\ref{amsa}.


\end{pf}





Для завершения доказательства леммы \ref{lstat} выберем подпоследовательность


$s_i, i=1,2,\dotsc$, последовательности $s=1,2,\dotsc$, так, что


\begin{gather}\label{lcond1}


I_j(u^i)\leq q^k_j+(q^{k,i}_j-q^k_j)+\gamma^{i,s_i},  


\ \ (q^{k,i}_j-q^k_j)+\gamma^{i,s_i}\to 0,\ \  j=1,\dotsc,l_k,\\


%


\gamma^{i,s_i}\to 0,\ \  \varepsilon^{i,s_i}\to


0,\ \  \alpha^i\gamma^{i,s_i}\to 0, 


\ \ I_0(u^{i,s_i})-\beta_k(q^k)\to 0,\ \  i\to\infty,\notag


\end{gather}


где введено обозначение $u^i\equiv u^{i,s_i}$. Одновременно без


ограничения общности в силу (\ref{lcond1}),


(\ref{**}), (\ref{*****}) можем считать, что


\begin{equation}\label{lcond2}


\mu_0^{i,s_i}\to 1,\ \  \alpha^i\mu_j^{i,s_i}\to\mu_j,  


\ \ j=1,\dotsc,l_k,


\end{equation}


где $\mu_j$ --- некоторые числа. Заметим, что числа


$\alpha^i$, $\mu_j^{i,s_i}$, $\mu_j$, вообще говоря, зависят от $k$. Умножая


соотношения~(\ref{***})--(\ref{*****}) леммы \ref{amsa} на $\alpha^i$,


пользуясь соотношениями (\ref{lcond1}) и (\ref{lcond2}),


можем переформулировать лемму \ref{amsa} в форме леммы \ref{lstat}.$\quad\square$








\begin{defns}


Последовательность управлений $u^i\in {\cal D}$, $i=1,2,\dotsc$, назовем


стационарной в задаче $(P^k_{q^k})$, если найдутся последовательность


неотрицательных чисел $\gamma^i\to0$, $i\to\infty$, и ограниченная


последовательность векторов $\mu^i\in R^{l_k}$, $i=1,2,\dotsc$,


$\mu^i\equiv(\mu^i_0,\mu^i_1,\dotsc,\mu^i_{l_k})$, такие, что $u^i\in {\cal


D}^{k,\gamma^i}_{q^k}$, выполнены соотношения (\ref{muic}),


(\ref{maxc}), где


$\lambda^i\equiv\sum\limits_{j=1}^{l_k}\mu^i_j\delta_{(x_{k,j},y_{k,j})}$, и


все предельные точки последовательности $\mu^i$, $i=1,2,\dotsc$, отличны от


нуля.


\end{defns}





Введем множества множителей


$L^{k, \nu}_{q^k}\equiv\Big\{-\sum\limits_{j=1}^{l_k}\mu_je^j\in R^{l_k}: 


\mu=(\mu_0,\mu_1,\dotsc,\mu_{l_l})\in R^{l_k+1}$, $\mu\neq 0$,  


$\mu_0=\nu$, существует стационарная в задаче $(P^k_{q^k})$ последовательность,


для которой соответствующая ей, согласно определению стационарной 


последовательности, последовательность векторов $\mu^i$, $i=1,2,\dotsc$,


имеет вектор $\mu$ в качестве своей предельной точки$\Big\}$, $\nu=0,1$;


$M^{k, 0}_{q^k}\equiv L^{k, 0}_{q^k}\cup\{0\}$,


$M^{k, 1}_{q^k}\equiv L^{k, 1}_{q^k}$. 


Непосредственно из леммы~\ref{lstat}, определения обобщенных 


субдифференциалов в смысле~\cite{mordukhovich88} и произвола в выборе


последовательностей $r^{k,i}$, $\sigma^{k,i}$, $i=1,2,\dotsc$, следует








\begin{thms}\label{subd}


Пусть $\beta_k(q^k)<+\infty$. Тогда


$\partial\beta_k(q^k)=\partial\beta_k(q^k)\cap M^{k, 1}_{q^k}$,


$\partial^\infty\beta_k(q^k)=\partial^\infty\beta_k(q^k)\cap M^{k, 0}_{q^k}$,


где $\partial\beta_k(q^k)$ и $\partial^\infty\beta_k(q^k)$ являются


соответственно обычным и сингулярным обобщенным субдифференциалами 


в точке $q^k$ $\cite{mordukhovich88}$.


\end{thms}





Из данной теоремы и теоремы 2.1 в \cite{mordukhovich88} вытекает








\begin{cors}\label{lip}


Если в некоторой окрестности $O_{q^k}$ точки $q^k$ все задачи $(P^k_{y^k})$,


$y^k\in O_{q^k}$, нормальны, т.\,е. $M^{k, 0}_{y^k}=\{0\}$, $y^k\in O_{q^k}$,


причем множества $M^{k, 1}_{y^k}$ равномерно ограничены постоянной $K$ 


в некоторой норме $\|\cdot\|$ (напр., евклидовой $|\cdot|$), то 


функция значений $\beta_k$ является липшицевой в сопряженной норме на


$O_{q^k}$ с той же постоянной $K$.


\end{cors}





Следствием теоремы~\ref{subd}, как и в~\cite{sumin00c}, является следующая теорема,


доказательство которой может быть проведено в полном соответствии со схемой


доказательства аналогичного результата в (\cite{sumin00c}, \S\,4).








\begin{thms}


Если задача $(P_q)$ нормальна, то ее функция


значений $\beta$ липшицева в окрестности точки $q\in C(X)$.


\end{thms}





Сформулируем, наконец, два достаточных условия нормальности в задаче $(P_q)$


(напр., \cite{sumin00b}).








\begin{thms}


Пусть в задаче $(P_q)$ правая часть $f$ имеет вид


$f(x,y,z,p,r,v)\equiv f_1(x,y)z+f_2(x,y)p+f_3(x,y)r+f_4(x,y,v)$,


а функция $G$ выпукла по~$z$. Пусть, кроме того, выполнено одно из 


следующих условий\/${:}$ $1)$~функция $\Phi$ выпукла по $z$, и 


найдется такое управление $u_*\in {\cal D}$,


что $\Phi(x,y,z[u_*](x,y))<q(x,y)$ \; $\forall$ $(x,y)\in X;$


$2)$~функция $\Phi$ имеет вид $\Phi(x,y,z)\equiv\Phi_1(x,y)z+\Phi_2(x,y)$


и в задаче $(P_q)$


найдется нестационарная последовательность $u^i\in {\cal D}^{\gamma^i}_q$,


$i=1,2,\dotsc$, $\gamma^i\to0$, $i\to\infty$. Тогда задача $(P_q)$ нормальна.


\end{thms}








\section{Иллюстративный пример}





Пусть задана оптимизационная задача


\begin{gather*}


I_0(u)\equiv z_2[u](1, 1)\to\inf,\ \  I_1(u)(x, y)\equiv-z_1[u](x, y)\leq


0,\ \  (x, y)\in \Pi\equiv[0,1]\times[0,1];\\


%


z_{1 xy}=(z_{1 x}+1)u(x, y),


\ \ z_{2 xy}=z_1-u^2(x, y),\ \  z(0, y)=z(x, 0)=0;\\


%


z=(z_1, z_2),\ \ u(x, y)\in U\equiv[-1, +1],


\ \  (x,y)\in\Pi\equiv[0, 1]\times[0, 1].


\end{gather*}


Можно показать, что в этой задаче не существует измеримого по Лебегу 


оптимального управления и что значение задачи равно $-1$.





Принцип максимума для м.\,п.\,р. теоремы \ref{Pontryagin} здесь будет иметь вид


\begin{multline*}


\int_\Pi\max\limits_{v\in[-1, +1]}\{(\mu^i\eta_{01}[u^i](x, y)+


\eta_{11}[u^i, \lambda^i](x, y))[(v-u^i(x, y))z_{1 x}[u^i](x, y)+\\


%


+(v-u^i(x, y))]+(\mu^i\eta_{02}[u^i](x, y)+\eta_{12}[u^i, \lambda^i](x, y))


((u^i(x, y))^2-v^2)\} dx\, dy\leq\gamma^i,


\end{multline*}


где $\gamma^i\to 0$, $i\to\infty$, а сопряженные функции


$\eta_0[u^i]\equiv(\eta_{01}[u^i], \eta_{02}[u^i])$,


$\eta_1[u^i,\lambda^i]\equiv(\eta_{11}[u^i,\lambda^i], \eta_{12}[u^i,\lambda^i])$


--- решения сопряженных систем


\begin{gather*}


\eta_1(x, y)-\int^1_yu^i(x, \xi_2)\eta_1(x, \xi_2) d\xi_2-


\int_x^1\int^1_y\eta_2(\xi_1, \xi_2) d\xi_1\,d\xi_2=0,\ \ \eta_2(x, y)=-1;\\


%


\eta_1(x, y)-\int_y^1u^i(x, \xi_2)\eta_1(x, \xi_2) d\xi_2-


\int_x^1\int^1_y\eta_2(\xi_1, \xi_2) d\xi_1\,d\xi_2=


\int^1_x\int^1_y\lambda^i(d\xi_1\,d\xi_2),\\


%


\eta_2(x, y)=0,


\end{gather*}


соответственно. Эти соотношения могут быть переписаны в виде


\begin{gather*}


\int_\Pi\max\limits_{v\in[-1, +1]}\{(\mu^i\eta_{01}[u^i](x, y)+


\eta_{11}[u^i, \lambda^i](x, y))[(v-u^i(x, y))z_{1 x}[u^i](x, y)+\\


%


+(v-u^i(x, y))]-\mu^i((u^i(x, y))^2-v^2)\} dx\, dy\leq\gamma^i,\\


%


\eta_{01}[u^i](x, y)-\int_y^1u^i(x, \xi_2)\eta_{01}(x, \xi_2) d\xi_2=


-(1-x)(1-y),\\


%


\eta_{11}[u^i](x, y)-\int_y^1u^i(x, \xi_2)\eta_{11}(x, \xi_2) d\xi_2=


\int^1_x\int^1_y\lambda^i(d\xi_1\,d\xi_2),\\


%


\mu^i+|\lambda^i|=1,\ \  \mes\{(x, y)\in\Pi : z_1[u^i](x, y)


\geq-\gamma^i\}\to0,\ \  i\to\infty,


\end{gather*}


где мера $\lambda^i$ сосредоточена на множестве


$\Pi^i\equiv\{(x, y)\in\Pi : |z_1[u^i](x, y)|\leq\gamma^i\}$,  $\gamma^i\to 0$,


$i\to\infty$.





Элементарный анализ этих соотношений показывает, что последовательность


$u^i\in {\cal D}$,\linebreak $i=1,2,\dotsc$, будет заведомо им удовлетворять, если


\begin{gather}\label{somecond}


\mu^i\eta_{01}[u^i](x, y)+\eta_{11}[u^i, \lambda^i](x, y)\to 0


\text{ по мере на } \Pi, \\


%


|u^i(x, y)|\to 1\text{ по мере на } \Pi,\ \  \mu^i+|\lambda^i|=1,\notag\\


%


\mes\{(x, y)\in\Pi : z_1[u^i](x, y)\geq-\gamma^i\}\to 0,\ \  i\to\infty,  


\ \ \supp\lambda^i\subset\Pi^i.\notag


\end{gather}


Одной из подходящих последовательностей пар $(\mu^i, \lambda^i)$ будет


последовательность, удовлетворяющая соотношениям


$$


\mu^i\to 1/2,\ \  |\lambda^i|\to 1/2,


\ \  \lambda^i[\{(\wt{x}, \wt{y})\in\Pi : x\leq\wt{x}\leq 1,  


\ y\leq\wt{y}\leq 1\}]\to (1/2)(1-x)(1-y).


$$


Для нее, очевидно, выполняется первое из соотношений (\ref{somecond}). При


этом, т.\,к. последовательность $u^i$, $i=1,2,\dotsc$, должна удовлетворять и


остальным соотношениям (\ref{somecond}), в этом случае с необходимостью


должно быть $|z_1[u^i]|^{(0)}_\Pi\to 0$, $i\to\infty$. Поскольку


$$


z_1[u](x, y)=\int^x_0\bigg(\exp\bigg(\int_0^yu(\xi_1,\xi_2)


d\xi_2\bigg)-1\bigg) d\xi_1,


$$


то одной из последовательностей, удовлетворяющей всем требованиям, является,


например, последовательность


$$


u^i(x,y)\equiv\begin{cases}


+1, & y\in(\frac{j-1}{2i}, \frac{j}{2i}),\ x\in[0, 1],\ j=1,3,\dotsc,2i-1;\\


-1, & y\in(\frac{j-1}{2i}, \frac{j}{2i}),\ x\in[0, 1],


\ j=2,4,\dotsc,2i;\ i=1,2,\dotsc


\end{cases}


$$


В то же время последовательность


$$


v^i(x,y)\equiv\begin{cases}


+1,& x\in(\frac{j-1}{2i},\frac{j}{2i}),\ y\in[0, 1],\ j=1,3,\dotsc,2i-1;\\


-1,& x\in(\frac{j-1}{2i},\frac{j}{2i}),\ y\in[0, 1],


\ j=2,4,\dotsc,2i;\  i=1,2,\dotsc,


\end{cases}


$$


последним требованиям не удовлетворяет, т.к.


$|z_1[v^i]|^{(0)}_\Pi\not\to 0$, $i\to\infty$. При этом


обе указанные последовательности в слабой норме $|\cdot|_w$


(\cite{warga}) сходятся к одному и тому же обобщенному управлению


$\nu(x, y)\equiv(\delta_{-1}+\delta_{+1})/2$. Итак, в данном


примере разные последовательности управлений сходятся к одному и тому же


обобщенному управлению, но одна из них является м.\,п.\,р. и идентифицируется с


помощью принципа максимума, а другая таковой не является.








\begin{thebibliography}{99}





\bibitem{gam77}


Гамкрелидзе Р.В. {\it Основы оптимального управления.}


-- Тбилиси: Изд-во Тбилисск. ун-та, 1977. -- 256~с.


\bibitem{warga}


Варга Дж. {\it Оптимальное управление


дифференциальными и функциональными уравнениями.} -- М.: Наука, 1977.  -- 624~с.


\bibitem{young}


Янг Л. {\it Лекции по вариационному исчислению и теории оптимального


управления.} -- М.: Мир, 1974.  -- 488~с.


\bibitem {SuminD}


Сумин М.И. {\it Математическая теория субоптимального управления


распределенными системами}\/: Дисс.~$\dotsc$ докт. физ.-матем. наук.


Нижний Новгород: Нижегородский гос. ун-т, 2000.  -- 350~с.


\bibitem{danmatv}


Данилова О.А., Матвеев А.С. {\it Нетрадиционные условия существования


оптимального управления для системы Гурса--Дарбу} // Изв. РАН.


Сер. матем. -- 1998. -- Т.\,62. -- \No\,5. -- C.\,79--102.


\bibitem{sumin00b}


Сумин М.И. {\it Субоптимальное управление полулинейными эллиптическими


уравнениями с фазовыми ограничениями.~$\mathrm I$. Принцип максимума для


минимизирующих последовательностей, нормальность} // Изв. вузов. Математика.


-- 2000. -- \No\,6. -- C.\,33--44.


\bibitem{sumin00c}


Сумин М.И. {\it Субоптимальное управление полулинейными эллиптическими


уравнениями с фазовыми ограничениями.~$\mathrm{II}$. Чувствительность, типичность


регулярного принципа максимума} // Изв. вузов. Математика. -- 2000. -- \No\,8.


-- C.\,52--63.


\bibitem{plotnikovsumin82} 


Плотников В.И., Сумин М.И. {\it О построении


минимизирующих последовательностей в задачах управления системами с


распределенными параметрами} // Журн. вычисл. матем. и матем. физ. -- 1982.


-- Т.\,22. -- \No\,1. -- С.\,49--56.


\bibitem{sumin97a} 


Сумин М.И. {\it Субоптимальное управление системами с


распределенными параметрами: минимизирующие последовательности, функция


значений} // Журн. вычисл. матем. и матем. физ. -- 1997. -- Т.\,37. -- \No\,1. --


С.\,23--41.


\bibitem{sumin97b} 


Сумин М.И. {\it Субоптимальное управление


системами с распределенными параметрами: свойства нормальности, субградиентный


двойственный метод} // Журн. вычисл. матем. и матем.


физ. -- 1997. -- Т.\,37. -- \No\,2. -- C.\,162--178.


\bibitem {altifom} 


Алексеев В.М., Тихомиров В.М., Фомин С.В.


{\it Оптимальное управление.} -- М.: Наука, 1979. -- 432~с.


\bibitem{clarke} 


Кларк Ф. {\it Оптимизация и негладкий анализ.} -- М.: Наука, 1988. -- 280~с.


\bibitem{mordukhovich88} 


Мордухович Б.Ш.


{\it Методы аппроксимаций в задачах оптимизации и управления.} -- М.: 


Наука, 1988. -- 360~с.


\bibitem{mordukh}


Mordukhovich B.S., Shao Y. {\it Nonsmooth sequential 


analysis in Asplund spaces} // Trans. Amer. Math. Soc. -- 1996. -- V.\,346.


-- \No\,4. -- P.\,1235--1280.


\bibitem{viplotnikovvisumin} 


Плотников В.И., Сумин В.И. {\it Оптимизация объектов


с распределенными параметрами, описываемых системами Гурса--Дарбу} // Журн.


вычисл. матем. и матем. физ. -- 1972. -- Т.\,12. -- \No\,1. -- С.\,61--77.


\bibitem{irkutyane}


{\it Условия экстремума и конструктивные методы решения в


задачах оптимизации гиперболических систем} / По ред. О.В.\,Васильева. --


Новосибирск: Наука, 1993. -- 196~с.


\bibitem{tihom}


Тихомиров В.М. {\it Выпуклый анализ} // Итоги науки и техн.


ВИНИТИ. Современ. пробл. матем. Фундаментальные


направления. -- 1987. -- Т.\,14. -- С.\,5--101.


\bibitem{plotsumv72}


Плотников В.И., Сумин


В.И. {\it Проблемы устойчивости нелинейных систем Гурса--Дарбу} // Дифференц.


уравнения. -- 1972. -- Т.\,8. -- \No\,5. -- С.\,845--856.


\bibitem{sumin1986}


Сумин М.И. {\it О минимизирующих последовательностях в задачах


оптимального управления при ограниченных фазовых координатах} // Дифференц.


уравнения. -- 1986. -- Т.\,22. -- \No\,10. -- С.\,1719--1731.


\bibitem{ekeland}


Ekeland I. {\it On the variational principle} // J. Math.


Anal. Appl. -- 1974.-- V.\,47. -- \No\,2. -- P.\,324--353.


\bibitem{sumqual} 


Сумин М.И. {\it Оптимальное управление объектами, описываемыми квазилинейными


эллиптическими уравнениями} // Дифференц. уравнения. -- 1989. -- Т.\,25.


-- \No\,8. -- С.\,1406--1416.





\end{thebibliography}





\vskip 2cm





\post{\it Нижегородский государственный}{\it Поступила}


\post{\it университет}{15.10.2002}





\label{end}


\end{document}








